% ------------------------------------------------------------------
% Diagrama de transicoes do AFD do validador de URLs (M_URL)
% Requer: \usetikzlibrary{automata, positioning, arrows.meta, fit}
% ------------------------------------------------------------------
\begin{figure}[htbp]
  \centering
  \resizebox{0.95\textwidth}{!}{%
  \begin{tikzpicture}[>={Stealth[round]}, shorten >=1pt, auto,
      every state/.style={minimum size=0.95cm, inner sep=1pt},
      initial text={},
      every edge/.style={draw, ->},
      every node/.style={font=\small}]

    % --- Linha 1: reconhecimento de "http" / "https" ---
    \node[state, initial] (q0) at (0,0)   {$q_0$};
    \node[state]          (q1) at (2.4,0) {$q_1$};
    \node[state]          (q2) at (4.8,0) {$q_2$};
    \node[state]          (q3) at (7.2,0) {$q_3$};
    \node[state]          (q4) at (9.6,0) {$q_4$};
    \node[state]          (q5) at (12,0)  {$q_5$};

    % Agrupamento: estados de prefixo (q1..q5)
    \node[draw, dashed, rounded corners, fit=(q1)(q5), inner sep=8pt,
          label={[font=\footnotesize]above:estados de prefixo $q_1,\dots,q_5$}] (pref) {};

    % --- Linha 2: separador "://" e rotulo em andamento ---
    \node[state] (q9) at (2.4,-3.4) {$q_9$};
    \node[state] (q8) at (4.8,-3.4) {$q_8$};
    \node[state] (q7) at (7.2,-3.4) {$q_7$};
    \node[state] (q6) at (9.6,-3.4) {$q_6$};

    % --- Linha 3: apos o ponto (extensao) ---
    \node[state] (q10) at (2.4,-6.8) {$q_{10}$};
    \node[state] (q11) at (4.8,-6.8) {$q_{11}$};
    \node[state, accepting] (q12) at (7.2,-6.8) {$q_{12}$};

    % Agrupamento: estados de extensao (q11, q12)
    \node[draw, dashed, rounded corners, fit=(q11)(q12), inner sep=8pt,
          label={[font=\footnotesize]above:letras ap\'os o ponto}] (ext) {};

    % Prefixo
    \path (q0) edge node {\texttt{h}} (q1)
          (q1) edge node {\texttt{t}} (q2)
          (q2) edge node {\texttt{t}} (q3)
          (q3) edge node {\texttt{p}} (q4)
          (q4) edge node {\texttt{s}} (q5);

    % Separador "://"
    \path (q4) edge node[left]  {\texttt{:}} (q6)
          (q5) edge node[right] {\texttt{:}} (q6)
          (q6) edge node[above] {\texttt{/}} (q7)
          (q7) edge node[above] {\texttt{/}} (q8)
          (q8) edge node[above] {$A$} (q9);

    % Rotulo em andamento (q9)
    \path (q0) edge node[below left, pos=0.45] {demais $A$} (q9);
    \draw[->] (pref.south -| q1) -- node[right, pos=0.35] {demais $W$} (q9);
    \path (q9) edge[loop left] node {$W$} (q9);

    % Ponto -> q10
    \path (q9)  edge[bend left=22]  node[right] {\texttt{.}} (q10)
          (q10) edge[bend left=22]  node[left]  {$D$} (q9);
    \draw[->] (pref.south west) to[out=-100, in=180, looseness=1.2]
          node[left, pos=0.5] {\texttt{.}} (q10.west);

    % Extensao
    \path (q10) edge node[above] {$L$} (q11)
          (q11) edge node[above] {$L$} (q12)
          (q12) edge[loop right] node {$L$} (q12);
    \draw[->] (ext.south) to[bend left=35] node[below] {\texttt{.}} (q10.south);
    \draw[->] (ext.north west) -- node[right, pos=0.45] {$D$, \texttt{-}} (q9);
  \end{tikzpicture}}
  \caption{AFD $M_{URL}$ do validador de URLs. $L$ = letras, $D$ = d\'igitos,
  $A = L \cup D$ e $W = A \cup \{\texttt{-}\}$. ``Demais $A$'' e ``demais $W$''
  indicam os s\'imbolos que n\~ao t\^em seta pr\'opria saindo daquele estado
  (por exemplo, de $q_1$ n\~ao se repete o \texttt{t}). Transi\c{c}\~oes n\~ao
  desenhadas levam ao estado de erro $q_e$.}
  \label{fig:afd-url}
\end{figure}